% from ega0-23-fr.tex
\oldpage[0\textsubscript{IV-1}]{219}

\begin{proof}
On applique (23.2.4) en prenant $B=\widehat A$~; comme les
$B/\mathfrak q_i$ sont des anneaux locaux noethériens complets, ce sont des
anneaux japonais (23.1.5) et $B$ est un $A$-module fidèlement plat
(0\textsubscript{I}, 7.3.5)~; en outre on a
$C_\lambda\otimes_AB=\widehat C_\lambda$ (Bourbaki, \emph{Alg. comm.},
chap.~III, §~3, n\textsuperscript{o}~4, th.~3 et chap.~IV, §~2,
n\textsuperscript{o}~5, cor.~3 de la prop.~9).
\end{proof}

\begin{corollary}[23.2.6]
\label{0.23.2.6-fr}
Sous les hypothèses de (23.2.5), la fermeture intégrale $A'$ de $A$ dans
$K'$ est un anneau semi-local~; si $\mathfrak m$ est l'idéal maximal de $A$,
alors, pour tout idéal maximal $\mathfrak m'$ de $A'$,
$A'/\mathfrak m'$ est une extension finie de $A/\mathfrak m$.
\end{corollary}

\begin{proof}
Le fait que l'homomorphisme
$(\widehat C_\alpha)_{\mathrm{red}}\to
(\widehat C_\lambda)_{\mathrm{red}}$ soit bijectif pour
$\lambda\geqslant\alpha$ entraîne que le nombre des idéaux maximaux de
$C_\lambda$ est constant pour $\lambda\geqslant\alpha$ et que si
$\mathfrak m_\alpha$ est un idéal maximal de $C_\alpha$ et
$\mathfrak m_\lambda$ l'unique idéal maximal de $C_\lambda$ au-dessus de
$\mathfrak m_\alpha$, les corps $C_\alpha/\mathfrak m_\alpha$ et
$C_\lambda/\mathfrak m_\lambda$ sont canoniquement isomorphes. La conclusion
résulte de ce que $\mathfrak m'=\varinjlim\mathfrak m_\lambda$ et
$A'/\mathfrak m'=\varinjlim(C_\lambda/\mathfrak m_\lambda)$
(0\textsubscript{III}, 10.3.1.1.3).
\end{proof}

\begin{proposition}[23.2.7]
\label{0.23.2.7-fr}
\textnormal{(Y. Mori).} Soient $A$ un anneau local noethérien intègre, $K$ son
corps des fractions, $A'$ la clôture intégrale de $A$. Alors $A'$ est un
anneau de Krull (Bourbaki, \emph{Alg. comm.}, chap.~VII, §~1) semi-local~;
autrement dit, il existe une famille $(V_\lambda)$ d'anneaux de valuation
discrète ayant $K$ pour corps des fractions, tels que
$A'=\bigcap_\lambda V_\lambda$ et que tout $x\in K$ appartienne à tous les
$V_\lambda$ sauf pour un nombre fini d'entre eux. En outre, il existe un
sous-anneau $C$ de $K$ qui est une $A$-algèbre finie, et telle que les
$V_\lambda$ soient les clôtures intégrales des anneaux
$C_{\mathfrak p_\lambda}$, où $(\mathfrak p_\lambda)$ est la famille des
idéaux premiers de hauteur $1$ de l'anneau local $C$.
\end{proposition}

Prouvons d'abord le lemme suivant~:

\begin{lemma}[23.2.7.1]
\label{0.23.2.7.1-fr}
Soient $A$, $B$ deux anneaux, $\varphi:A\to B$ un homomorphisme faisant de $B$
un $A$-module fidèlement plat. On suppose que $A$ est intègre~; alors, si
$C=B_{\mathrm{red}}$, l'homomorphisme composé
$\psi:A\to B\to B_{\mathrm{red}}=C$ est injectif et se prolonge en un
homomorphisme injectif du corps des fractions $K$ de $A$ dans l'anneau total
des fractions $R$ de $C$~; en outre, si $C'$ est la fermeture intégrale de
$C$ dans $R$, $C'\cap K$ est la clôture intégrale de $A$.
\end{lemma}

\begin{proof}
Comme $A$ est intègre et $\varphi$ injectif, l'intersection de $\varphi(A)$
et du nilradical $\mathfrak N$ de $B$ est réduite à $0$, et comme
$C=B/\mathfrak N$, $\psi$ est injectif. Tout $a\neq0$ dans $A$ étant non
diviseur de zéro dans $A$ ne l'est pas non plus dans $B$ par platitude
(0\textsubscript{I}, 6.3.4)~; on en déduit que $a$ n'est pas non plus diviseur
de $0$ dans $C$, car si l'on avait $ax\in\mathfrak N$ pour un
$x\notin\mathfrak N$ dans $B$, on en tirerait $a^nx^n=0$ pour un entier $n$,
ce qui contredit ce qui précède puisque $x^n\neq0$. On peut donc prolonger
$\psi$ en un homomorphisme injectif de $K$ dans $R$. Pour prouver la dernière
assertion, notons qu'il est clair que la clôture intégrale $A'$ de $A$ est
contenue dans $C'\cap K$. Inversement, soit $x\in C'\cap K$~; $x$ est donc
entier sur $C$, et \emph{a fortiori} sur $B$, autrement dit $B[x]$ est une
$B$-algèbre finie~; d'ailleurs, $K$ s'identifie à un sous-anneau de
$B\otimes_AK$, et le sous-anneau $B[x]$ de $B\otimes_AK$ s'identifie à
$B\otimes_AA[x]$ par platitude~; on en conclut que $A[x]$ est un $A$-module
de type fini (Bourbaki, \emph{Alg. comm.}, chap.~I\textsuperscript{er}, §~3,
n\textsuperscript{o}~6, prop.~11), donc que $x\in A'$.
\end{proof}

Ce lemme étant établi, nous allons l'appliquer à l'injection canonique de $A$
dans son complété $\widehat A$, qui est un $A$-module fidèlement plat~;
$B=(\widehat A)_{\mathrm{red}}=\widehat A/\mathfrak N$ (où $\mathfrak N$ est
le nilradical de $\widehat A$) est un anneau local noethérien complet réduit,
dont l'anneau total des

